% Responsable: P3  (versión completa en extendido/04_lexer_concordancia_arboles.tex)
\section{Léxico, concordancia y visualización}
\label{sec:lexer_concordancia_arboles}

\paragraph{Lexer.} Es el de la Fase~1 adaptado: el alfabeto incluye \emph{á, é, í, ó, ú, ü, ñ},
el texto se normaliza a Unicode NFC, y la categoría ya no sale de una expresión regular sino del
léxico, porque la misma forma puede ser sustantivo, verbo o adjetivo. Una palabra desconocida
no es un error léxico: recibe la categoría \texttt{X} y el parser rechaza la oración.

\paragraph{Léxico.} \texttt{data/lexicon.json} tiene 221 palabras con su categoría y sus rasgos
(género, número, persona y tiempo): 52 sustantivos, 58 adjetivos, 7 verbos en presente y
pretérito (70 formas) y las palabras funcionales. Cinco palabras son ambiguas, como \emph{bajo}.

\paragraph{Analizador semántico (concordancia).} \texttt{spanish/agreement.py} es la versión de
la Fase~2 del analizador semántico (\texttt{semantico.py}) y la tercera etapa del pipeline. Así
como en la Fase~1 la tabla de símbolos verificaba reglas que la CFG no expresa, este analizador
recorre el árbol y verifica que, en cada sintagma nominal, determinante, sustantivo y adjetivo coincidan en género y número, y que
el núcleo del sujeto coincida con el verbo en número y persona. Un rasgo ausente concuerda con
todo (adjetivos invariables como \emph{grande}, posesivos sin género). Como trabaja sobre el
árbol, resuelve la concordancia a distancia: acepta \emph{El perro de los niños come} y rechaza
\emph{*El perro de los niños comen}.

\paragraph{Visualización.} \texttt{spanish/trees.py} dibuja el árbol en la terminal
(\texttt{nltk}) o como imagen con Graphviz, con un color por categoría y las palabras que no
concuerdan en rojo
(figura~\ref{fig:arbol_concordancia}).

\begin{figure}[h]
  \centering
  \includegraphics[width=0.42\linewidth]{arbol_concordancia}
  \caption{Árbol generado por el pipeline para \emph{*La elefante grises duerme}: se deriva,
    pero el analizador semántico marca en rojo las dos palabras que no concuerdan.}
  \label{fig:arbol_concordancia}
\end{figure}
